%=======================================================================
%  CHAPTER 8 -- CONCLUSION
%=======================================================================
\chapter{Conclusion and Future Work}
\label{ch:conclusion}

\section{Summary of the Work}
\label{sec:conc-summary}

This thesis constructed and evaluated a uniform causal anomaly detection pipeline to determine whether, and more usefully when, causal structure provides a measurable advantage over standard machine learning. The evaluation was executed as a combinatorial grid: three causal discovery algorithms (PCMCI, GES, CAM-UV) crossed with three regression families (linear, polynomial, RBF) across three domains (Robotics, Healthcare, Industrial). So that differences would be attributable to the algorithms, not to tuning, the scoring mechanism, the thresholding rule and the attribution depth were held constant, enforced programmatically by an automated audit on every build.

\section{Answers to the Research Questions}
\label{sec:conc-rqs}

\subsection*{RQ1 --- Algorithm effectiveness and ranking stability}
There is no universally optimal causal algorithm. PCMCI takes the highest score on Robotics (0.8752) and Industrial (0.9151), while GES takes the highest on Healthcare (0.8114). The regression-variant ranking within an algorithm is statistically established: GES linear separates from GES polynomial on Healthcare at $p = 0.001119$, but the cross-algorithm ranking within a domain is not testable at all, because the subsampling policy means the algorithms do not score the same observations. The domain ordering is empirical; the algorithmic ranking is not established.

\subsection*{RQ2 --- Interaction between discovery and regression}
The regression family interacts with both the discovery algorithm and the domain. Across the nine algorithm--domain cells the RBF variant takes five, linear two and polynomial two, and GES prefers a different variant in each of the three domains. Neither axis dominates: which regression family is best cannot be stated without naming both the algorithm and the domain.

\subsection*{RQ3 --- Causal versus non-causal detection}
The advantage is conditional on the character of the faults. Causal methods outperform the best tuned neural baseline on Healthcare, 0.8114 against 0.7077. On Industrial the causal pipeline is beaten by the neural baselines and cannot be separated from the parameter-free trivial control, 0.9151 against 0.9286, which lies inside its bootstrap interval. On Robotics the two are not separable.

\subsection*{RQ4 --- Attribution quality}
The causal pipeline names the variable whose modelled relationship broke instead of the one that deviated most. It returns \texttt{TabletTouchX} in 9 of 9 Robotics configurations for the LED attack and \texttt{KneePicthTemp} in 8 of 9 for the joint attack, and neither is the variable that moved most: the tablet channel ranks 78th of 249 by standardised deviation and the knee temperature 4th.

The more striking half of this result is that the attribution does not require the graph to contain the relation that would explain it. Six of the nine LED configurations return the correct variable from graphs holding no LED channel whatsoever, the variance-ranked cap applied by GES and CAM-UV on Robotics having excluded every one. PCMCI, which applies no cap and sees all 244 variables, returns the same answer. A correct explanation was produced from a structure that does not contain the explanation.

Those are two faults of the forty-three this work scores, and the answer does not hold at that scale. Scored against the documented cause of each of the twenty-two anomaly classes that carry one, across all three domains, the intrinsic attribution is no more relevant than post-hoc attribution over the neural baselines. Paired over classes, only the electrocardiogram carries enough scorable classes to decide the question, and there the intrinsic attribution is measurably behind under PCMCI and indistinguishable from post-hoc attribution under GES and CAM-UV. So the answer to RQ4 is yes on the question as asked, since the attribution does name the variable whose relationship broke and not the one that moved most, and no on the question a reader will ask next, since in the four comparisons where both sides rank a pool of the same size the causal pipeline never names it correctly more often than the post-hoc one. What it does retain is set out in Section~\ref{sec:attribution} and returned to in the closing remarks.

\subsection*{Why CAM-UV places as it does}
This is not one of the four questions, but it is what explains the answer to the first of them. CAM-UV never takes the highest score in any domain, and its latent confounder mechanism shows three distinct regimes that account for it. On Robotics it saturates, assigning 80.7\% of all variable pairs to the latent set and leaving no observed-parent edges at all, so the structural constraint that is supposed to do the work is absent. On Industrial it starves, with 0.67 samples per variable leaving the independence test no power to flag anything. Healthcare falls between them, saturating at 59.1\% and reaching 0.7861. The synthetic ablation confirms the reading without excusing it: recovery of a planted confounder is monotone in sample size, never succeeding at 200 and succeeding at every significance level tested by 1500.

\section{Contributions Revisited}
\label{sec:conc-contributions}

Section~\ref{sec:contributions} claimed three, and each is revisited here against what the evaluation actually delivered rather than what it set out to.

C1, the systematic evaluation, holds. The 27 configurations differ only along the two
declared methodological axes, across domains of 244, 52 and 12 variables, and Section~\ref{sec:method-overview} describes the audit that enforces this on every build rather than asserting it. The linearity assumption was relaxed in turn and not left in place, and the answer to RQ2 is that no regression family dominates: which one performs best cannot be stated without naming both the algorithm and the domain.

C2, the accuracy against explainability tradeoff, holds in a more interesting form than it was claimed. The detection half is conditional and the condition is measurable in advance, which is the finding this work would be cited for: the advantage of causal structure over the best tuned neural baseline shrinks monotonically as the score of a parameter-free control rises, so a question usually settled by building both methods can be estimated before either is built. The explainability half does not track it. Attribution succeeds in 9 of 9 configurations for the LED fault, six of those from graphs containing no LED channel at all, and in 8 of 9 for the joint fault, including on domains where the detection advantage has gone. The two halves of the tradeoff therefore come apart, and a practitioner choosing on detection alone would discard the explanatory capability for nothing. That capability is not a higher score: scored across the twenty-two anomaly classes that carry a documented cause, and over pools of matched size, the intrinsic attribution is no more relevant than the post-hoc one, and what it offers instead is a named broken relationship and, on the electrocardiogram, agreement that is specific to the record being explained. One qualification is recorded in Section~\ref{sec:disc-xai}: the joint-fault half depends on the ranking metric this domain uses and would have read as a failure under the other.

C3, the comparison against neural architectures, holds on both axes it was claimed on. On detection the three tuned autoencoders and the parameter-free control were evaluated under the identical label-blind rule, and the ordering between families changes with the domain and is not fixed. On explanation, post-hoc SHAP over those baselines returned the variable that moved furthest for the joint fault, naming RShoulderPitch, the single largest raw deviation in the system, in all three architectures, where the causal attribution named the thermal cause ranked fourth. That single fault is the clearest case for the distinction the comparison exists to draw, and it does not generalise. Scored against the documented cause of each of the twenty-two anomaly classes that carry one, and restricted to the comparisons in which both sides rank a pool of the same size, the post-hoc attribution is nowhere less relevant than the intrinsic one. The contribution therefore holds as a comparison rather than as a demonstration of causal superiority.

Two further findings emerged instead of being sought, and neither was claimed as a contribution. Downstream detection performance cannot validate an upstream causal graph: a structurally empty graph produced competitive detection, and two graphs built on incompatible premises produced indistinguishable scores, so a pipeline reporting only an AUC-ROC offers no evidence that its recovered structure is physically meaningful. And treating the subsampling rate as a property of each algorithm forecloses cross-algorithm significance testing, quantified in Table~\ref{tab:delong-summary}: 54 of the 108 within-domain comparisons admit no paired test, split identically 18 against 18 in every domain, and every unpairable one involves PCMCI. The empirical bounding of CAM-UV on real data, with its measured saturation and starvation regimes and the sample-size floor below which recovery of a planted confounder never succeeded, is reported in Appendix~\ref{app:confounders} and not carried here.

\section{Relation to Prior Work}
\label{sec:disc-prior}

This thesis does not establish that causal anomaly detection outperforms or underperforms deep learning in general, and no such claim is made. What it establishes is narrower and, for a reader deciding how to evaluate a method, more useful: that the comparison as conventionally constructed omits the control that determines how its result should be read.

The Industrial domain demonstrates this directly. Reported in the standard form, best causal configuration against best tuned neural baseline, the result is 0.9151 against 0.9997, and the conclusion available from those two numbers is that causal methods underperform on industrial process data. Adding a parameter-free control changes the conclusion rather than qualifying it. The trivial mean absolute $z$-score reaches 0.9286, above the causal pipeline and requiring no learned parameters whatsoever. The finding is therefore not that the causal method failed but that neither method's structure did work a threshold could not: the faults are amplitude excursions, and there is little relational signal for either to exploit. Those are different conclusions with different consequences for a practitioner, and only the second is supported once the control is present.

This is consistent with three separate critiques of evaluation practice, which are worth keeping distinct because they identify different faults. Wu and Keogh~\cite{wu2023flawed} argue that the benchmarks themselves manufacture an appearance of progress. Schmidl et al.~\cite{schmidl2022anomaly}, re-implementing 71 algorithms across 976 datasets, report no clear winner between algorithm families. Kim et al.~\cite{kim2022rigorous} show that the point-adjustment protocol inflates scores to the point where a random score looks competitive. None of the three establishes that elaborate detectors lose to trivial baselines as a general rule; what they establish jointly is that the protocol, not the architecture, frequently decides the reported ordering. The contribution here is to demonstrate, on one domain, that causal pipelines are not exempt from that finding, a possibility the causal discovery literature has had little occasion to test, since parameter-free controls are rarely reported alongside them.

A second point concerns attribution. Evaluation of explanations frequently rests on whether an explanation appears reasonable to a reader, a practice criticised by Doshi-Velez and Kim~\cite{doshivelez2017rigorous} and shown by Adebayo et al.~\cite{adebayo2018sanity} to admit explanations that fail elementary sanity checks. The documented ground truth used here avoids that dependence. It also yields a result that should give pause: as Section~\ref{sec:disc-structure} reports, the pipeline attributes the LED attack correctly in all nine configurations, including six whose graphs contain no LED channel at all and therefore cannot represent the relation that would explain the attribution. A correct explanation was produced from a structure that does not contain the explanation. This supports the caution urged by Pham et al.~\cite{pham2024howfar} in the root-cause analysis literature: neither a strong detection score nor a plausible-looking graph is evidence that the recovered structure is physically accurate.

The scope of the claim should be stated exactly. Nothing here shows that published comparisons reached incorrect conclusions on their own data. What is shown is that a comparison constructed the way those comparisons are constructed would have reached a different conclusion on this data, and that the difference is attributable to a single omitted control that carries no learned parameters at all.

\section{Limitations}
\label{sec:conc-limitations}

Six limitations bound what the preceding answers are worth, and they are not of equal weight.

The heaviest concerns the inputs rather than the pipeline. Scoring resolution differs by algorithm: PCMCI subsamples at 74, 7 and 1 where GES and CAM-UV are fixed at 10. Within a domain, therefore, the three algorithms are not scoring the same series, which forecloses the paired testing RQ1
would otherwise rest on. A second asymmetry of the same kind is the ridge penalty, which varies by regression variant on the Robotics domain; it is identical across the three algorithms inside every cell, so RQ1 is untouched, but a within-algorithm comparison on that domain changes two things at once.

Reproducibility is the next constraint. Each causal configuration was run once, and CAM-UV is not exactly reproducible between runs, so the interval estimates come from resampling the scores of a single run instead of from repeated runs. Two of the 27 configurations cannot be distinguished from chance at all: GES with linear features on Robotics, 0.5906 with an interval of $[0.4873, 0.6761]$, and PCMCI with polynomial features on Healthcare, 0.5075 with an interval of $[0.4429, 0.5663]$.

The remaining two concern how far the findings travel. There is one dataset per domain, the robotics faults are injected and the industrial faults are simulated, though the healthcare classes are real recorded diagnoses carrying a clinician's label. And the Healthcare result, which is the strongest in the work, depends substantially on exact algebraic identities among the ECG leads. Those identities are a favourable case. Noisy field data will not generally supply them.

\section{Future Work}
\label{sec:conc-future}

Three lines follow directly from what was found. In the near term the comparison should be re-run at a single common scoring resolution on the domains where that is computationally feasible, which would remove the primary threat to internal validity rather than merely bounding it. In the medium term the CAM-UV latent set should be thresholded or ranked by confidence instead of integrating every flagged pair, a change the Robotics saturation result directly motivates.

The most consequential is a fourth domain, selected for a known low trivial-control score. That is the direct test of the conditional hypothesis in Section~\ref{sec:disc-when}: if the causal pipeline reliably wins where amplitude excursions are absent, the hypothesis survives a test it could have failed. Proposing the experiment that would refute the claim is the appropriate way to leave it.

\section{Closing Remarks}
\label{sec:conc-closing}

On detection, the question this evaluation answers is not whether causal anomaly detection is superior to deep learning, but when. The answer is legible in a control that costs nothing to compute. Where faults present as amplitude excursions, structure is unnecessary and a threshold suffices. Where faults are relational, the variables remaining within their normal ranges while the dependencies between them break, structure provides an advantage that neither a threshold nor a reconstruction model recovered. Establishing the condition is more useful than claiming the victory.

On explanation the answer is less comfortable, and not the one this work expected to give. Scored against the documented cause of each of the twenty-two anomaly classes that carry one, and compared only where both sides rank a pool of the same size, the intrinsic causal attribution was nowhere more relevant than post-hoc attribution over the neural baselines. The case for reading an explanation off the structure the detector already uses cannot therefore rest on it being right more often, because on this evidence it is not. Nor does the evidence support the opposite case with any confidence. On the one domain carrying enough scorable classes to test, the gap separates from zero only under PCMCI; on the other two the question cannot be put at all, since three attack classes and five documented actuators cannot reach the 5\% level however the numbers fall. What the three benchmarks jointly supply is three attack classes, five documented actuators and fourteen localisable diagnoses, and that is not enough ground truth to rank two families of explanation against each other. The contribution this work can defend is the measurement showing as much, together with the property the causal attribution retains regardless: it names a relationship rather than a magnitude, and it does so without fitting a second model to approximate the first.

What the causal side retains is narrower and worth having. On the electrocardiogram its configurations agree with one another about the specific record in front of them, carrying a margin of 0.267 between agreement within a class and agreement any class would produce, where the neural architectures carry 0.047 and return something close to a fixed ranking of the leads. And the explanation it gives is of a different kind: a named relationship that stopped holding, available without fitting a second model to approximate the first. That is a claim about what an explanation is, not about what it scores, and it is the one this thesis can support.
